\unnumchapter{Заключение}

\todo{Заключение~--- черновик, требует ручной доработки.}

В работе разработана система управления размещением нагрузки в федерации
Kubernetes-кластеров с учётом метрик надёжности. Система построена как
сквозной контур управления, связывающий наблюдаемость, показатели уровня
обслуживания, политики размещения, выбор кластера и аварийное
восстановление, и экспериментально оценена на воспроизводимых отказах. Все
целевые показатели достигнуты, а в центральном сценарии планировщик,
учитывающий текущие метрики кластеров, даёт на 45.6\% меньше ошибок, чем
наивный выбор.

Получены ответы на поставленные исследовательские вопросы.

Локальные показатели кластеров корректно агрегируются через Thanos и служат
входными данными планировщика. Однако глобальную доступность федерации так
вычислять нельзя: сумма серверных показателей не видит отказа кластера, и
доступность необходимо измерять в точке входа трафика.

Оба исследованных механизма дают измеримый эффект: переключение трафика
снижает потери запросов при отказе кластера на 87.1\%, а размещение с учётом
метрик снижает долю ошибок при деградации сети на 45.6\%. Встроенное
самовосстановление Kubernetes с проверками по умолчанию, напротив, усиливает
сетевой отказ.

Цикл аварийного восстановления выполняется программно и воспроизводим: время
обнаружения отказа не превышает 13~с, время восстановления~--- 43~с во всех
прогонах. Времена отсчитываются от событий инструмента внесения отказа, а
фазы инцидента записываются в каталог.

Политики размещения выражены требованиями нагрузки и описанием кластера и
исполняются на двух независимых рубежах; ни один нарушитель не был размещён.

Планировщик, использующий метрики из каталога, превосходит выбор по очереди
при тех же обязательных ограничениях в каждой паре прогонов. Его
преимущество устойчиво к выбору весов, пока оценка учитывает метрику,
чувствительную к типу отказа.

Направления дальнейшей работы: исследование зависимости эффекта
планировщика и накладных расходов контура от числа кластеров; перенос уже
размещённых нагрузок при деградации кластера; иная оценка кластеров без
трафика вместо нейтральных значений; оценка доверительных интервалов при
большем числе повторов; резервирование управляющего контура и шлюза.
